\documentclass[18pt]{article}
\usepackage[utf8]{inputenc}
\usepackage[T1]{fontenc}
\usepackage{ragged2e}
\usepackage{caladea}
\usepackage{graphicx}
\usepackage{longtable}
\usepackage{wrapfig}
\usepackage{rotating}
\usepackage{epigraph}
\usepackage[normalem]{ulem}
\usepackage{hyperref}
\usepackage{amsmath}
\usepackage{amssymb}
\usepackage{capt-of}
\usepackage{hyperref}
\usepackage{fancyhdr}

\title{
 \includegraphics[scale=.55, trim={10cm, 0, 10cm, 0}]{./assets/imagem.jpg}
  \par
   Novena a Santo Expedito}
  \date{ Data de Início: 07/04 Data Litúrgica: 16/04}

% Comando para fazer "Sumário" não aparecer no Sumário.
\renewcommand{\contentsname}{Sumário}
\begin{document}

\thispagestyle{empty} %zera a primeira página

\pagestyle{fancy}
\fancyhf{} % clear existing header/footer entries
\fancyfoot[LO, CE]{
\includegraphics[scale=0.2]{./assets/cross.png} Santa Bernadette Soubirous, rogai por nós!
}
% Place Page X of Y on the right-hand
% side of the footer
\fancyfoot[R]{\thepage}

\newpage

\tableofcontents

\centering
\vfill
Visite-nos no Telegram: \url{https://t.me/CotidieNovena}
\newpage

\newpage


%%%%%%%%%%%%%%%%%%%%%%%%%%%%%%%%%%%%% História %%%%%%%%%%%%%%%%%%%%%%%%%%%%%%%%%%%%%%%%%%%

\begin{justify}

 \begin{center}
  \section{História}\label{sec:História} % (fold)
 \end{center}

A XII Legião tinha o nome de Fulminata, nome que significa algo como Vem como um raio. A Fulminata tinha cerca de 8 mil homens contando soldados, escravos e cavaleiros. No tempo de Expedito, ela defendia as fronteiras orientais do Império Romano contra os bárbaros asiáticos. Expedito comandou a XII Legião de 296 a 303 d. C. Para ser comandante de uma Legião Romana era preciso muita competência e bravura. Tanto que, alguns anos antes, a mesma XII Legião tinha sido comandada por um Imperador romano, Marco Aurélio, numa campanha onde hoje é a Eslováquia. Sabe-se que Expedito era um líder competente. Seu cargo equivaleria hoje ao de um general. Ele se tornou famoso por manter a disciplina dos soldados e todos o respeitavam. Por outro lado, como a maioria dos soldados romanos, o Comandante Expedito tinha uma vida devassa, rodeada de luxo, prazeres e fama.

O primeiro contato do Comandante Expedito com o Cristianismo aconteceu dentro da própria XII Legião. Com efeito, uma parte dos soldados da XII era formada de cristãos. Além disso, em suas andanças pelas fronteiras orientais do império, Expedito teve ainda mais contato com o cristianismo. E, para completar, a XII Legião teve um soldado chamado Polieucto de Melitene, que morreu mártir no ano 193. A semente do Cristianismo e do sangue dos mártires nunca cai na terra em vão.

Expedito era um líder competente na condução da XII Legião tanto nos tempos de paz quanto  nas batalhas. Um comandante vitorioso na carreira militar. Porém, quanto à sua vida espiritual, tinha o vício da procrastinação, isto é, deixar para depois, adiar. Ele simpatizava com a mensagem de Jesus. Admirava os ensinamentos do Mestre de Nazaré e via no Evangelho palavras que ninguém jamais tinha dito antes na história humana. Por isso, ele pensava em um dia converter-se de verdade. Esse dia, porém, ficava sempre para mais tarde, era sempre adiado.

Depois de alguns anos procrastinando, Expedito foi tocado pela graça de Deus. Certa noite teve um sonho que mudou sua vida. No sonho, com um corvo representando o espírito do mal, grasnava diante dele a palavra cras, do latim, que significa amanhã, deixe sua conversão para amanhã. O corvo grasnava forte e parecia poderoso. Porém, de repente, Expedito decidiu e pisoteou o corvo dizendo: hodie, que significa hoje, em latim. O Comandante Expedito acordou do sonho decidido e confirmou sua conversão. Por isso ele é considerado o Santo das causas urgentes. Convertido, ele continuou por um tempo ainda chefe da sua legião, conseguindo converter seus soldados também ao cristianismo.

Com a conversão de Expedito e da sua tropa, o imperador Diocleciano começou a perseguir o Santo e seus soldados. A importância de seu posto fazia dele uma influência muito forte a favor do Cristianismo dentro do Império Romano. Por isso, ele se tornou alvo especial do Imperador.

Santo Expedito foi preso pela ordem de Diocleciano e foi forçado a renunciar à sua nova fé. Porém, ele não renunciou. Seus castigos começaram pela flagelação romana: 39 chicotadas com o flagrus, chicote que dilacera a pele e causa hemorragia. Expedito tinha aplicado este mesmo castigo a bandidos e indisciplinados. Agora, ele os recebia por causa de Jesus Cristo. E ele permaneceu firme e julgando-se indigno de sofrer o mesmo castigo que Jesus sofrera, aplicado por soldados romanos, como estava acontecendo com ele. Por fim, não renunciando à sua fé, Santo Expedito foi decapitado com espada, por ordem do Imperador Diocleciano, no dia 19 de abril de 303, em Melitene na Armênia.

Como Santo Expedito foi morto por causa de Jesus Cristo, ele se tornou um mártir da Igreja e santo reconhecido oficialmente. Sua bravura diante dos sofrimentos por causa da fé serviu de exemplo para grande parte dos soldados de sua Legião, fazendo-os permanecer firmes em sua fé. O exemplo de Santo Expedito arrastou milhares de cristãos na Armênia e, logo, ele passou a ser venerado como santo, o santo das causas urgentes.

Santo Expedito é representado como um Comandante Romano, vestindo uma túnica branca, uma armadura de superior e um manto vermelho sobre os ombros. Em sua mão direita ele levanta uma cruz com a palavra Hodie (Hoje). Na mão esquerda ele tem uma palma, representando o martírio e a vitória dos mártires. Seu pé direito pisa sobre um corvo, que grita a palavra Cras, (amanhã). A imagem simboliza a grande mensagem de Santo Expedito: Não adie sua conversão, não deixe para amanhã aquilo que deve ser feito hoje, não procrastine!

\vfill

\begin{center}
 \href{https://santo.cancaonova.com/santo/santas-perpetua-e-felicidade/}{Fonte: Canção Nova}
\end{center}

%%%%%%%%%%%%%%%%%%%%%%%%%%%%%%%%%%%%% Orações %%%%%%%%%%%%%%%%%%%%%%%%%%%%%%%%%%%%%%%%%%%

\newpage
\begin{center}
 \section{Orações}\label{sec:Orações} % (fold)
\textit{Em nome do Pai, e do Filho, e do Espírito Santo. Amém.}
\end{center}

\subsection{Oração Incial}\label{sec:Oração_Inicial} % (fold)

Ó Minha Senhora e minha Mãe, Vós manifestastes em Lourdes a grandeza de vosso poder e a imensidade de vossa bondade, e por isso todos os homens vão a Vós e Vós os curais. Não permitais que eu me separe de Vós. Mãe, é o que Vos peço mais do que tudo, do fundo de minha alma. Pois, dessa forma, não me afastarei de Nosso Senhor Jesus Cristo, vosso Filho, para o qual Vós sois, por vontade d’Ele, canal necessário, medianeira universal e onipotência suplicante. Nossa Senhora de Lourdes, rogai por nós! Santa Bernadette, rogai por nós!


\subsection{Ato de Contrição}\label{sec:Ato_de_Contrição} % (fold)

Meu Pai e Meu Senhor Jesus Cristo, caridade sem fim, eu sinceramente me arrependo dos meus pecados. Concedei-me, portanto, o perdão dos meus pecados e a graça que peço pelos méritos das dores de Sua Mãe amorosa e pelas virtudes de Seu mártir Santo Expedito.


\subsection{Oração Final}\label{sec:Oração_Final} % (fold)


Ó minha Mãe e minha Senhora de Lourdes, sei que eu merecia ser abandonado. Mas sei que Vós existis e que pedis para o homem não aquilo a que ele tem direito, mas aquilo a que ele não tem direito. Vós pedis para o homem o perdão a que ele não tem direito, a generosidade a que ele perdeu o direito, o afago a que seus atos de virtude não dão título. Pois bem, tudo isso Vós obtendes para eles, pelos méritos de vosso sorriso, junto ao vosso Divino Filho. E eu sei que em determinado momento sairei disto e continuarei a ir para frente. Ó minha Mãe, vede onde me deixei cair! Minha Mãe, não conseguirei nada enquanto Vós não me ajudardes. Ajudai-me, ajudai-me! Amém.

Lembre-se, graciosa Virgem Maria, que nunca se ouviu dizer que qualquer um que se refugiou sob Vossa proteção, implorou a Vossa ajuda e procurou Vossa intercessão foi deixado desamparado. Inspirado por essa confiança, eu recorro a Vós, Virgem das virgens, minha Mãe. Para Vós eu venho, diante de Vós estou, pecador e triste. Mãe do Verbo Encarnado, não desprezeis as minhas palavras, mas graciosamente ouvi-me e atendei a minha oração. Amém.

\textbf{\textit{Rezar três Pai-Nossos em honra a Santíssima Trindade}}

\textbf{Em nome do Pai, e do Filho, e do Espírito Santo}

\vfill

\subsection{Orações de Cada Dia}



\subsection*{Primeiro Dia}

\textbf{\textit{\nameref{sec:Oração_Inicial}}} % (fold)

Oh, glorioso mártir Santo Expedito! Pela fé viva que foi concedida por Deus, peço-lhe para despertar a mesma fé no meu coração, para que eu também acredite, sinceramente, que há Deus, mas muito especialmente para que eu possa ser salvo de pecar contra Ele.

\textbf{\textit{\nameref{sec:Ato_de_Contrição} e \nameref{sec:Oração_Final}}} % (fold)

\subsection*{Segundo Dia}

\textbf{\textit{\nameref{sec:Oração_Inicial}}} % (fold)

Oh, glorioso mártir Santo Expedito! Pela esperança dada por Deus, orai para que aqueles de pouca fé possam ser modificados por alguns raios de esperança, para que eles também recebem as coisas eternas. Por favor, orai para que ardente esperança em Deus seja também me dada e me segure firme no meio de sofrimentos.

\textbf{\textit{\nameref{sec:Ato_de_Contrição} e \nameref{sec:Oração_Final}}} % (fold)

\subsection*{Terceiro Dia}

\textbf{\textit{\nameref{sec:Oração_Inicial}}} % (fold)

Oh, glorioso mártir Santo Expedito! Pelo amor infinito que Nosso Senhor plantou em seu coração, por favor, retire do meu todos os grilhões atados pelo mundanismo, que sem eles eu possa amar somente a Deus em toda a eternidade. Amém.

\textbf{\textit{\nameref{sec:Ato_de_Contrição} e \nameref{sec:Oração_Final}}} % (fold)

\subsection*{Quarto Dia}

\textbf{\textit{\nameref{sec:Oração_Inicial}}} % (fold)

Oh, glorioso mártir Santo Expedito, que sabia perfeitamente bem o ensinamento do Mestre Divino para carregar a Cruz e segui-Lo! Peça-Lhe as graças de que eu preciso para que possa lutar contra as minhas próprias paixões.

\textbf{\textit{\nameref{sec:Ato_de_Contrição} e \nameref{sec:Oração_Final}}} % (fold)

\subsection*{Quinto Dia}

\textbf{\textit{\nameref{sec:Oração_Inicial}}} % (fold)

Oh, glorioso mártir Santo Expedito, pelas graças abundantes que recebeste do céu, que possa me conservar todas as suas virtudes, concedei-me também que possa ser livre de todos os sentimentos que bloqueiam o meu caminho para o Céu.

\textbf{\textit{\nameref{sec:Ato_de_Contrição} e \nameref{sec:Oração_Final}}} % (fold)

\subsection*{Sexto Dia}

\textbf{\textit{\nameref{sec:Oração_Inicial}}} % (fold)

Oh, glorioso mártir Santo Expedito! Pelos sofrimentos e humilhações que recebeste para o amor de Deus, concedei-me esta graça também que é muito agradável a Deus e me liberta da raiva e da dureza de coração, que é a pedra de tropeço da minha alma.

\textbf{\textit{\nameref{sec:Ato_de_Contrição} e \nameref{sec:Oração_Final}}} % (fold)

\subsection*{Sétimo Dia}

\textbf{\textit{\nameref{sec:Oração_Inicial}}} % (fold)

Oh, Santo Expedito, meu protetor! Eu coloco a minha esperança no fato de que minhas petições podem ser concedidas se forem para o meu próprio bem. Por favor, pedi ao Senhor, por intercessão da Virgem Santíssima, o perdão dos meus pecados e a graça de mudar a minha vida, particularmente a graça: (mencionar aqui a graça particular desejada). Eu prometo seguir Seus exemplos e propagar esta devoção.

\textbf{\textit{\nameref{sec:Ato_de_Contrição} e \nameref{sec:Oração_Final}}} % (fold)

\subsection*{Oitavo Dia}

\textbf{\textit{\nameref{sec:Oração_Inicial}}} % (fold)

Oh, glorioso mártir Santo Expedito! Por meio de um desejo puro, que reinou em todos os seus sentimentos, palavras e ações, por favor, guiai-me na minha busca incessante para a glória de Deus e o bem dos meus semelhantes.

\textbf{\textit{\nameref{sec:Ato_de_Contrição} e \nameref{sec:Oração_Final}}} % (fold)

\subsection*{Nono Dia}

\textbf{\textit{\nameref{sec:Oração_Inicial}}} % (fold)

Oh, glorioso mártir Santo Expedito, que foi muito amado pela Rainha do Céu, que nada lhe foi negado, pedi a Ela, por favor, oh meu advogado, que, pelos sofrimentos de seu Filho Divino e Suas próprias tristezas, eu possa receber, neste dia, a graça que peço, mas acima de tudo, a graça de morrer sem cometer algum pecado mortal. Amém.

\textbf{\textit{\nameref{sec:Ato_de_Contrição} e \nameref{sec:Oração_Final}}} % (fold)



\centering

\vfill
\subsection*{Créditos:}
\href{https://pocketterco.com.br/index.php/terco/novena-a-santo-expedito-inicia-em-10-de-abril}{Pocket Terço}


\end{justify}

\end{document}
